\section{Analysis of Model~C$_2$ ($\theta_1>0,\,\gamma_1=\gamma_2=\theta_2=0$)}

In this section, we will derive $P(x,y)$ for this model where only true abandonments coming from Class-$1$ take place, while no abandonments in Class-$2$ or jockeying of any sort are taken into account.

\begin{figure}[H]
    \centering
    \resizebox{0.62\textwidth}{!}{%
        % Model C₂: γ₁=γ₂=θ₂=0, θ₁>0
% Class-1 customers abandon the queue at rate θ₁ (dotted, upward from Queue 1).
% No jockeying. Class-2 never abandons.
\begin{tikzpicture}[>=Latex, thick]

    % ── Queue 1 (top, Class-1, priority) ──────────────────────────────────────
    \draw (0, 1.9) -- (2, 1.9) -- (2, 1.1) -- (0, 1.1);
    \foreach \x in {0.2, 0.4, ..., 1.8} { \draw[thin] (\x, 1.8) -- (\x, 1.2); }
    \draw[->] (-1.5, 1.5) -- (-0.1, 1.5) node[midway, above] {$\lambda_1$};
    % θ₁: Class-1 abandonment (dotted, upward)
    \draw[->, dotted] (1, 2.0) -- (1, 3.0) node[above] {$\theta_1$};

    % ── Queue 2 (bottom, Class-2) ──────────────────────────────────────────────
    \draw (0, -1.1) -- (2, -1.1) -- (2, -1.9) -- (0, -1.9);
    \foreach \x in {0.2, 0.4, ..., 1.8} { \draw[thin] (\x, -1.2) -- (\x, -1.8); }
    \draw[->] (-1.5, -1.5) -- (-0.1, -1.5) node[midway, above] {$\lambda_2$};
    % No θ₂: Class-2 does not abandon

    % ── Server ────────────────────────────────────────────────────────────────
    \node[draw, circle, minimum size=1.2cm] (server) at (5, 0) {$\mu$};
    \draw[->] (2.2,  1.5) -- (server);
    \draw[->] (2.2, -1.5) -- (server);
    \draw[->] (server) -- (7, 0);

\end{tikzpicture}
%
    }
    \caption{Queue layout for Model~C$_2$. Class-1 customers abandon the system
    at rate $\theta_1$ per waiting customer (dotted arrow, upward from Queue~1).
    Class-2 never abandons ($\theta_2=0$) and there is no jockeying.}
    \label{fig:model-c2-queue}
\end{figure} The rationale behind using this combination of parameters is that, thanks to the fact that Class-$2$ still sees Poisson arrivals, the same sort of probabilistic arguments we used for Model A will hold. Moreover, we will see that the resulting equation when plugging this combination of parameters in our fundamental equation (Equation~\eqref{eq:C2:fundamental}) yields a \textit{first-order linear ODE} in $x$ and, $P_y(y)$ can be taken out of the integral, saving us from having to solve an equation of the Volterra family.

\begin{theorem}
    Consider a queueing system with $2$ priority classes with Poisson arrival rates $\lambda_1$ and $\lambda_2$, where class-$1$ customers have non-preemptive priority over those of class-$2$, exponentially distributed service times of rate $\mu$ independently of the class, and per-customer abandonment of class-$1$ customers from the queue at rate $\theta_1$ (with no class-$2$ abandonment and no jockeying). Under the stability condition
    \begin{equation}
        \rho_2\cdot{}_1F_1\!\left(1;\,\tfrac{\mu}{\theta_1}+1;\,\tfrac{\lambda_1}{\theta_1}\right) \;<\; 1
        \quad\Longleftrightarrow\quad
        \lambda_2\,\mathbb{E}[B_C] \;<\; 1,
    \end{equation}
    the probability generating function $P(x,y)$ of such a system is given by
    \begin{align}
        \begin{split}
            P(x,y) \;&=\; \frac{e^{\lambda_1 x/\theta_1}\,\mu\,\pi(0,0)\,(1-y)}{\theta_1\, x^{\mu/\theta_1}(1-x)^{\lambda_2(1-y)/\theta_1}\bigl(\widetilde{B}_C(\lambda_2(1-y))-y\bigr)}\\[4pt]
            &\quad\times \Bigl[\widetilde{B}_C(\lambda_2(1-y))\,\mathcal{I}_1(x,y) - \bigl(1-\widetilde{B}_C(\lambda_2(1-y))\bigr)\,\mathcal{I}_2(x,y)\Bigr],
        \end{split}
    \end{align}
    where
    \begin{align*}
        \mathcal{I}_1(x,y) &\;=\; \int_0^x e^{-\lambda_1 t/\theta_1}\,t^{\,\mu/\theta_1-1}(1-t)^{\,\lambda_2(1-y)/\theta_1}\,dt, \\
        \mathcal{I}_2(x,y) &\;=\; \int_0^x e^{-\lambda_1 t/\theta_1}\,t^{\,\mu/\theta_1}(1-t)^{\,\lambda_2(1-y)/\theta_1-1}\,dt,
    \end{align*}
    $\widetilde{B}_C(s)$ is the LST of the busy period of the $M/M/1+M$ subsystem of class-$1$ given by Equation~\eqref{eq:C2:Bc_LST}, and $\pi(0,0)$ is given by Corollary~\ref{cor:C2:pi00} below.
\end{theorem}

\begin{corollary}\label{cor:C2:pi00}
    The limiting probability $\pi(0,0)$ of the joint state $(N_1,N_2)=(0,0)$ with the server busy, and the empty-system probability $\pi_0$, are given by
    \begin{align}
        \pi(0,0) \;&=\; \frac{(\lambda_1+\lambda_2)\bigl(1 - \lambda_2\,\mathbb{E}[B_C]\bigr)}{\mu\,\bigl(1 + \lambda_1\,\mathbb{E}[B_C]\bigr)} \;=\; (\rho_1+\rho_2)\,\frac{1 - \lambda_2\,\mathbb{E}[B_C]}{1 + \lambda_1\,\mathbb{E}[B_C]}, \label{eq:C2:pi00}\\[4pt]
        \pi_0 \;&=\; \frac{1 - \lambda_2\,\mathbb{E}[B_C]}{1 + \lambda_1\,\mathbb{E}[B_C]}, \label{eq:C2:pi0}
    \end{align}
    where $\mathbb{E}[B_C]$ is the mean busy period of the $M/M/1+M$ subsystem of class-$1$, given by Equation~\eqref{eq:C2:Bc_mean}. In the limit $\theta_1\to 0^+$, $\mathbb{E}[B_C]\to (\mu-\lambda_1)^{-1}$ and the expressions above recover the classical Model~A results $\pi(0,0) = \rho(1-\rho)$ and $\pi_0 = 1-\rho$.
\end{corollary}

\subsection{Setup: reduction to a first-order linear ODE in $x$}

Plugging $\theta_1>0,\ \gamma_1=\gamma_2=\theta_2=0$ into the fundamental equation~\eqref{eq:gen:fundamental_equation} (both partial-derivative terms collapse onto the $\theta_1$ contribution) leaves
\begin{equation}
    \left[ (\lambda_1+\lambda_2+\mu)xy - \mu y - \lambda_1 x^2 y - \lambda_2 x y^2 \right] P(x,y)
    + \theta_1\, xy(x-1)\,\frac{\partial}{\partial x} P(x,y)
    \;=\; \mu (x-y)\,P_y(y) - \mu x(1-y)\,\pi(0,0),
    \label{eq:C2:fundamental}
\end{equation}
where the boundary function $P_y(y)=P(0,y)=\sum_{n_2=0}^{\infty} \pi(0,n_2)\,y^{n_2}$ is, as in Model A, the partial PGF restricted to states where the server is busy and the priority queue is empty ($N_1=0$).

For any fixed $y$, Equation~\eqref{eq:C2:fundamental} is a \emph{first-order linear ODE in $x$} whose unknowns are the function $P(x,y)$ on $[0,1]^2$ and the boundary trace $P_y(y)=P(0,y)$. The strategy in what follows is in two stages:
\begin{enumerate}
    \item Identify $P_y(y)$ — first by a probabilistic argument exploiting the $M/G/1$ structure seen by Class-$2$ (\S\ref{ssec:C2:prob}), then re-derived analytically from boundedness of $P(x,y)$ at $x=1$ (\S\ref{ssec:C2:anal}).
    \item Recover $P(x,y)$ by solving the ODE with $P_y(y)$ now known (\S\ref{ssec:C2:recovery}).
\end{enumerate}
We close the setup with the conditional-expectation decomposition that will drive the probabilistic derivation:
\begin{equation}
    P_y(y) \;=\; \mathbb{P}(\text{busy},\, N_1=0)\,\cdot\,\mathbb{E}\!\left[y^{N_2} \,\big|\, \text{busy},\, N_1=0\right].
    \label{eq:C2:Py_prob_decomp}
\end{equation}

\subsection{Probabilistic derivation of $P_y(y)$}\label{ssec:C2:prob}

The idea is, again, to adopt the perspective of a Class-$2$ customer waiting in line while the server is busy. Unlike in the case of jockeying, the customers from class-$1$ truly abandon the system, leaving the Poisson arrival process of Class-$2$ undisturbed. Therefore, customers in the non-priority class are driven by an $M/G/1$ process, with Poisson arrivals at rate $\lambda_2$ and general effective service times $B_C$ equivalent to the busy period of Class-$1$'s $M/M/1+M$ queue with arrival rate $\lambda_1$, service rate $\mu$ and per-customer abandonment rate $\theta_1$.

The LST and expectation of such busy period $B_C$ are known results:
\begin{align}
    \begin{split}
        \widetilde{B}_C(s) &= \cfrac{\mu}{\lambda_1+\mu+s-\cfrac{\lambda_1(\mu+\theta_1)}{\lambda_1+\mu+\theta_1+s-\cfrac{\lambda_1(\mu+2\theta_1)}{\lambda_1+\mu+2\theta_1+s-\dotsb}}},
    \label{eq:C2:Bc_LST}
    \end{split} \\
    \begin{split}
         \mathbb{E}[B_C] &= \sum_{k=1}^{\infty} \frac{\lambda_1^{\,k-1}}{\displaystyle\prod_{i=1}^{k}\bigl(\mu+(i-1)\theta_1\bigr)}
    \;=\;\frac{{}_1F_1(1;\,\mu/\theta_1+1;\,\lambda_1/\theta_1)}{\mu}.
    \label{eq:C2:Bc_mean}
    \end{split}
\end{align}
The second equality follows from expressing the product in terms of \textit{Kummer's confluent hypergeometric function of the first kind} ${}_1F_1$,
\begin{equation}
    {}_1F_1(a;b;z) \;=\; \sum_{n=0}^{\infty} \frac{a^{(n)}\,z^n}{b^{(n)}\,n!}, \qquad a^{(n)}\;=\;\prod_{k=0}^{n-1}(a+k) \quad(a^{(0)}=1),
\end{equation}
which provides a closed form for $\mathbb{E}[B_C]$.

The stability condition required for the Class-$2$ subsystem is that of the $M/G/1$ defined above, i.e.\ $\lambda_2\,\mathbb{E}[B_C]<1$:
\begin{equation}
    \rho_2 \cdot {}_1F_1\!\left(1;\,\frac{\mu}{\theta_1}+1;\,\frac{\lambda_1}{\theta_1}\right) \;<\; 1.
    \label{eq:C2:stability}
\end{equation}
Compared to Model A, this stability condition is strictly weaker than the classical $\rho_1+\rho_2<1$, which aligns with the fact that abandonments lower the load of the system: for all $\theta_1>0$, the series defining ${}_1F_1(1;\,\mu/\theta_1+1;\,\lambda_1/\theta_1)$ is termwise dominated by the geometric series for $1/(1-\rho_1)$, so ${}_1F_1(1;\,\mu/\theta_1+1;\,\lambda_1/\theta_1) < 1/(1-\rho_1)$. Abandonments in the priority queue therefore stabilise the non-priority queue.

With the LST and expectation of the effective service time seen by Class-$2$ customers, and thanks to the constant arrival rate $\lambda_2$ from the Poisson arrival process, we apply the Pollaczek--Khinchine formula:
\begin{equation}
    L_2(y) \;\equiv\; \mathbb{E}\!\left[y^{N_2} \,\big|\, \text{busy},\, N_1=0\right] \;=\; \frac{(1 - \lambda_2 \mathbb{E}[B_C])\,(1-y)\,\widetilde{B}_C(\lambda_2(1-y))}{\widetilde{B}_C(\lambda_2(1-y)) - y}.
    \label{eq:C2:pk_formula}
\end{equation}
Substituting~\eqref{eq:C2:pk_formula} into~\eqref{eq:C2:Py_prob_decomp} yields
\begin{equation}
    P_y(y) \;=\; \mathbb{P}(\text{busy},\,N_1=0)\,\left[\frac{(1-\lambda_2\mathbb{E}[B_C])\,(1-y)\,\widetilde{B}_C(\lambda_2(1-y))}{\widetilde{B}_C(\lambda_2(1-y)) - y}\right].
    \label{eq:C2:Py_unsimplified}
\end{equation}
By the definition of $P_y$, $P_y(0)=\pi(0,0)$. Evaluating~\eqref{eq:C2:Py_unsimplified} at $y=0$,
\begin{align}
    P_y(0) \;=\; \pi(0,0)
    &\;=\; \mathbb{P}(\text{busy},\,N_1=0)\,\left[\frac{(1-\lambda_2\mathbb{E}[B_C])\cdot 1 \cdot \widetilde{B}_C(\lambda_2)}{\widetilde{B}_C(\lambda_2)}\right] \nonumber \\
    &\;=\; \mathbb{P}(\text{busy},\,N_1=0)\,(1-\lambda_2\mathbb{E}[B_C]),
    \label{eq:C2:Py_at_zero}
\end{align}
so that $\mathbb{P}(\text{busy},\,N_1=0)=\pi(0,0)/(1-\lambda_2\mathbb{E}[B_C])$. Substituting this back into~\eqref{eq:C2:Py_unsimplified}, the $(1-\lambda_2\mathbb{E}[B_C])$ factor cancels and we are left with
\begin{equation}
    P_y(y) \;=\; \pi(0,0)\,\frac{(1-y)\,\widetilde{B}_C(\lambda_2(1-y))}{\widetilde{B}_C(\lambda_2(1-y)) - y}.
    \label{eq:C2:Py}
\end{equation}
This is structurally the same expression as in the probabilistic argument for Model~A (Equation~\eqref{eq:A:Py}), with the only difference being the effective service time: $B_C$ is now the busy period of an $M/M/1+M$ subsystem with abandonment rate $\theta_1$, rather than the busy period of a pure $M/M/1$.

\subsection{Analytical re-derivation of $P_y(y)$ via boundedness at $x=1$}\label{ssec:C2:anal}

For an analytic counterpart of the previous derivation, we put~\eqref{eq:C2:fundamental} in standard form:
\begin{align}
    \frac{\partial P}{\partial x} \;+\; \underbrace{\frac{(\lambda_1+\lambda_2+\mu)x - \mu - \lambda_1 x^2 - \lambda_2 xy}{\theta_1\,x(x-1)}}_{p(x;\,y)}\,P
    \;=\; \underbrace{\frac{\mu(x-y)\,P_y(y) - \mu x(1-y)\,\pi(0,0)}{\theta_1\,xy(x-1)}}_{q(x;\,y)}.
    \label{eq:C2:ODE_standard}
\end{align}
The numerator of $p(x;y)$ is the same polynomial we have seen in every model so far. With $x^*(y)$ and $x^+(y)$ as its roots (Equation~\eqref{eq:A:xstar}), we can write
\begin{equation}
    (\lambda_1+\lambda_2+\mu)x - \mu - \lambda_1 x^2 - \lambda_2 xy \;=\; -\lambda_1\bigl(x-x^*(y)\bigr)\bigl(x-x^+(y)\bigr),
\end{equation}
with $x^*(y)\,x^+(y) = \mu/\lambda_1$. Partial fraction decomposition gives
\begin{equation}
    \frac{-\lambda_1\bigl(x-x^*(y)\bigr)\bigl(x-x^+(y)\bigr)}{x(x-1)} \;=\; A + \frac{B}{x} + \frac{C}{x-1},
\end{equation}
where $A=-\lambda_1$ is read off from the ratio of leading coefficients of the quadratic terms in numerator and denominator, $B=\mu$ from setting $x=0$ and using $\lambda_1\,x^*x^+=\mu$, and $C=\lambda_2(1-y)$ from setting $x=1$ and using $(1-x^*)(1-x^+) = -\lambda_2(1-y)/\lambda_1$. Thus
\begin{equation}
    p(x;y) \;=\; \frac{1}{\theta_1}\!\left[-\lambda_1 + \frac{\mu}{x} + \frac{\lambda_2(1-y)}{x-1}\right].
    \label{eq:C2:pfd}
\end{equation}
Integrating term by term,
\begin{equation}
    \int p(x;y)\,dx \;=\; \frac{1}{\theta_1}\!\left[-\lambda_1 x + \mu\ln x + \lambda_2(1-y)\ln(1-x)\right],
\end{equation}
which yields the integrating factor
\begin{equation}
    \mu_I(x;y) \;=\; e^{-\lambda_1 x/\theta_1}\cdot x^{\alpha}\cdot (1-x)^{\beta(y)},
    \label{eq:C2:integrating_factor}
\end{equation}
where, in parallel with Model~$B_2$, we set
\begin{equation}
    \alpha \;=\; \frac{\mu}{\theta_1}, \qquad \beta(y) \;=\; \frac{\lambda_2(1-y)}{\theta_1}.
    \label{eq:C2:exponents}
\end{equation}
Under our stability assumptions both exponents are positive ($\alpha>0$ always; $\beta(y)\ge0$ for $y\in[0,1]$), so $\mu_I$ vanishes at both $x=0$ and $x=1$. From $\tfrac{d}{dx}[P\cdot\mu_I] = q\cdot\mu_I$, integration over $(0,x]$ reads
\begin{equation}
    P(x,y)\,\mu_I(x;y) - P(0,y)\,\mu_I(0;y) \;=\; \int_0^x q(t;y)\,\mu_I(t;y)\,dt,
\end{equation}
and the boundary term vanishes, yielding
\begin{equation}
    P(x,y) \;=\; \frac{1}{\mu_I(x;y)} \int_0^x q(t;y)\,\mu_I(t;y)\,dt.
    \label{eq:C2:Pxy_implicit}
\end{equation}

Now, in the same fashion as the root-cancellation argument used for Model A, we exploit boundedness of PGFs to determine $P_y(y)$ by observing the behaviour of $P(x,y)$ as $x\to 1$. Since $\mu_I(x;y)\to 0$, the only way to keep the right-hand side of~\eqref{eq:C2:Pxy_implicit} bounded is
\begin{equation}
    \int_0^1 q(t;y)\,\mu_I(t;y)\,dt \;=\; 0.
    \label{eq:C2:boundedness}
\end{equation}
Substituting the explicit forms of $q$ and $\mu_I$ into~\eqref{eq:C2:boundedness} and using $(1-t)^{b}/(t-1) = -(1-t)^{b-1}$ to simplify,
\begin{equation}
    \int_0^1 e^{-\lambda_1 t/\theta_1}\,t^{\,\alpha-1}(1-t)^{\,\beta(y)-1}\bigl[(t-y)\,P_y(y) - t(1-y)\,\pi(0,0)\bigr]\,dt \;=\; 0,
\end{equation}
with $\alpha$ and $\beta(y)$ as in~\eqref{eq:C2:exponents}. Introducing the auxiliary integrals
\begin{align}
    \mathcal{J}_0(y) &\;=\; \int_0^1 e^{-\lambda_1 t/\theta_1}\,t^{\,\alpha-1}(1-t)^{\,\beta(y)-1}\,dt, \\
    \mathcal{J}_1(y) &\;=\; \int_0^1 e^{-\lambda_1 t/\theta_1}\,t^{\,\alpha}(1-t)^{\,\beta(y)-1}\,dt,
\end{align}
and splitting the integrand, the boundedness condition~\eqref{eq:C2:boundedness} becomes
\begin{equation}
    P_y(y)\bigl[\mathcal{J}_1(y) - y\,\mathcal{J}_0(y)\bigr] \;=\; (1-y)\,\pi(0,0)\,\mathcal{J}_1(y),
\end{equation}
which gives
\begin{equation}
    P_y(y) \;=\; \pi(0,0)\,(1-y)\,\frac{\mathcal{J}_1(y)}{\mathcal{J}_1(y) - y\,\mathcal{J}_0(y)}.
    \label{eq:C2:Py_anal_raw}
\end{equation}

The continued-fraction representation of the busy period of an $M/M/1+M$ queue~\eqref{eq:C2:Bc_LST} admits the integral form
\begin{equation}
    \widetilde{B}_C(s) \;=\; \frac{\displaystyle\int_{0}^{1} t^{a}(1-t)^{b-1}e^{-ct}\,dt}{\displaystyle\int_{0}^{1} t^{a-1}(1-t)^{b-1}e^{-ct}\,dt}
    \;=\; \frac{B(a+1,b)\,{}_1F_1(a+1;\,a+1+b;\,-c)}{B(a,b)\,{}_1F_1(a;\,a+b;\,-c)},
    \label{eq:C2:Bc_Kummer}
\end{equation}
with $a=\alpha=\mu/\theta_1$, $b=s/\theta_1$, $c=\lambda_1/\theta_1$. Setting $s=\lambda_2(1-y)$ — so that $b=\beta(y)$ matches the exponent in~\eqref{eq:C2:exponents} — yields exactly the ratio of our $\mathcal{J}_i$:
\begin{equation}
    \widetilde{B}_C\bigl(\lambda_2(1-y)\bigr) \;=\; \frac{\mathcal{J}_1(y)}{\mathcal{J}_0(y)}.
\end{equation}
Dividing numerator and denominator of~\eqref{eq:C2:Py_anal_raw} by $\mathcal{J}_0(y)$ then recovers exactly Equation~\eqref{eq:C2:Py}. The boundedness route, which uses no probabilistic structure, reaches the same Pollaczek--Khinchine form only through the Beta-integral identity~\eqref{eq:C2:Bc_Kummer}, which is the analytic shadow of the $M/G/1$/busy-period decomposition used in \S\ref{ssec:C2:prob}.

\subsection{Recovery of $P(x,y)$ from $P_y(y)$}\label{ssec:C2:recovery}

With $P_y(y)$ now in hand, we return to the ODE~\eqref{eq:C2:fundamental} and substitute~\eqref{eq:C2:Py} into its right-hand side. Writing $\zeta(y) \equiv \widetilde{B}_C(\lambda_2(1-y))$ for brevity,
\begin{align*}
    \text{RHS}
    &\;=\; \mu(x-y)\,\pi(0,0)\,\frac{(1-y)\,\zeta(y)}{\zeta(y) - y} \;-\; \mu x(1-y)\,\pi(0,0) \\
    &\;=\; \mu\,\pi(0,0)\,(1-y)\!\left[\frac{(x-y)\,\zeta(y)}{\zeta(y) - y} - x\right] \\
    &\;=\; \mu\,\pi(0,0)\,(1-y)\,\frac{(x-y)\,\zeta(y) - x\bigl[\zeta(y) - y\bigr]}{\zeta(y) - y} \\
    &\;=\; \mu\,\pi(0,0)\,y(1-y)\,\frac{x - \zeta(y)}{\zeta(y) - y},
\end{align*}
which is structurally identical to the RHS found for Model A, though now with the more involved effective service time $B_C$. With $p(x;y)$ and $\mu_I(x;y)$ as in~\eqref{eq:C2:pfd} and~\eqref{eq:C2:integrating_factor}, the source $q(x;y)$ of the standard-form ODE now reads
\begin{equation}
    q(x;y) \;=\; \frac{\mu\,\pi(0,0)\,(1-y)\,\bigl[x-\zeta(y)\bigr]}{\theta_1\,x(x-1)\,\bigl[\zeta(y)-y\bigr]},
\end{equation}
and the implicit solution~\eqref{eq:C2:Pxy_implicit} becomes
\begin{equation}
    P(x,y) \;=\; \frac{1}{\mu_I(x;y)} \int_0^x q(t;y)\,\mu_I(t;y)\,dt.
\end{equation}

To evaluate the integral, we use the partial-fraction identity
\begin{align*}
    \frac{t-\zeta(y)}{t(t-1)}
    &\;=\; \frac{\zeta(y)(t-1) + t\bigl(1-\zeta(y)\bigr)}{t(t-1)}
    \;=\; \frac{\zeta(y)}{t} + \frac{1-\zeta(y)}{t-1},
\end{align*}
and define the auxiliary integrals
\begin{align}
    \mathcal{I}_1(x,y) &\;=\; \int_0^x e^{-\lambda_1 t/\theta_1}\,t^{\,\alpha-1}(1-t)^{\,\beta(y)}\,dt, \label{eq:C2:I1} \\
    \mathcal{I}_2(x,y) &\;=\; \int_0^x e^{-\lambda_1 t/\theta_1}\,t^{\,\alpha}(1-t)^{\,\beta(y)-1}\,dt, \label{eq:C2:I2}
\end{align}
with $\alpha$ and $\beta(y)$ as in~\eqref{eq:C2:exponents}. Since $\pi(0,0)$ and $\zeta(y)$ come out of the integral,
\begin{equation}
    \int_0^x q(t;y)\,\mu_I(t;y)\,dt \;=\; \frac{\mu\,\pi(0,0)\,(1-y)}{\theta_1\bigl(\zeta(y)-y\bigr)}\Bigl[\zeta(y)\,\mathcal{I}_1(x,y) \;-\; \bigl(1-\zeta(y)\bigr)\,\mathcal{I}_2(x,y)\Bigr].
\end{equation}
Dividing by $\mu_I(x;y)$ and restoring $\zeta(y)=\widetilde{B}_C(\lambda_2(1-y))$, we arrive at the closed form
\begin{align}
    \begin{split}
        P(x,y) \;=\; \frac{e^{\lambda_1 x/\theta_1}\,\mu\,\pi(0,0)\,(1-y)}{\theta_1\, x^{\alpha}(1-x)^{\beta(y)}\bigl(\widetilde{B}_C(\lambda_2(1-y))-y\bigr)}\\[4pt]
        \times \Bigl[\widetilde{B}_C(\lambda_2(1-y))\,\mathcal{I}_1(x,y) - \bigl(1-\widetilde{B}_C(\lambda_2(1-y))\bigr)\,\mathcal{I}_2(x,y)\Bigr],
    \end{split}\label{eq:C2:Pxy}
\end{align}
which is the statement of the theorem. This closed form is the product of a probabilistic derivation of $P_y(y)$ via Pollaczek--Khinchine, combined with an analytic integration of the ODE in $x$; the same result is reached from the purely analytic route of \S\ref{ssec:C2:anal}.

\begin{proof}[Proof of Corollary~\ref{cor:C2:pi00}]
    The boundary balance at the empty state $(N_1,N_2)=(0,0)$ with server idle, equating the rate-out (arrivals at total rate $\lambda_1+\lambda_2$) and the rate-in (service completion from $(0,0)$ with the server busy at rate $\mu$), reads
    \begin{equation}
        (\lambda_1+\lambda_2)\,\pi_0 \;=\; \mu\,\pi(0,0),
        \label{eq:C2:balance_zero}
    \end{equation}
    so $\pi(0,0) = (\rho_1+\rho_2)\,\pi_0$.

    To fix $\pi_0$, we exploit the $M/G/1$ structure used in \S\ref{ssec:C2:prob}. Class-$2$ customers form an $M/G/1$ with arrival rate $\lambda_2$ and service time $B_C$, whose utilisation is $\lambda_2\,\mathbb{E}[B_C]$. The event ``$M/G/1$ idle'' coincides with the absence of any Class-$2$ customer from the system, so
    \begin{equation}
        \mathbb{P}(\text{no Class-$2$ present}) \;=\; 1 - \lambda_2\,\mathbb{E}[B_C].
    \end{equation}
    Conditional on no Class-$2$ being present, the priority subsystem evolves freely as the same $M/M/1+M$ queue used to define $B_C$, with arrival rate $\lambda_1$, service rate $\mu$ and abandonment rate $\theta_1$. The stationary idle fraction of that subsystem is given by the standard renewal-theoretic ratio of mean idle period ($1/\lambda_1$) to mean cycle ($1/\lambda_1+\mathbb{E}[B_C]$):
    \begin{equation}
        \mathbb{P}(\text{Class-$1$ absent} \mid \text{no Class-$2$ present}) \;=\; \frac{1}{1+\lambda_1\,\mathbb{E}[B_C]}.
    \end{equation}
    The event $\{\text{server idle}, N_1=N_2=0\}$ is precisely the intersection of these two, so
    \begin{equation}
        \pi_0 \;=\; \bigl(1-\lambda_2\,\mathbb{E}[B_C]\bigr)\cdot\frac{1}{1+\lambda_1\,\mathbb{E}[B_C]} \;=\; \frac{1-\lambda_2\,\mathbb{E}[B_C]}{1+\lambda_1\,\mathbb{E}[B_C]},
    \end{equation}
    which is~\eqref{eq:C2:pi0}. Substituting into~\eqref{eq:C2:balance_zero} yields~\eqref{eq:C2:pi00}.

    For the Model A limit, $\theta_1\to 0^+$ gives $\mathbb{E}[B_C]\to(\mu-\lambda_1)^{-1}$, hence $\lambda_1\,\mathbb{E}[B_C]\to\rho_1/(1-\rho_1)$ and $1+\lambda_1\,\mathbb{E}[B_C]\to 1/(1-\rho_1)$; substituting into~\eqref{eq:C2:pi0} and~\eqref{eq:C2:pi00} recovers $\pi_0=1-\rho$ and $\pi(0,0)=\rho(1-\rho)$.
\end{proof}

\subsection{Mean queue lengths and mean waiting times}

\subsubsection*{Why the diagonal no longer gives a closed form}

Setting $x=y=z$ in the fundamental equation~\eqref{eq:C2:fundamental}, the
terms $\mu(x-y)P_y(y)$ and $\gamma_1 xy(x-y)\partial_x P$ both vanish, but
the abandonment convection $\theta_1 xy(x-1)\partial_x P$ does \emph{not}
(since $(x-1)|_{x=z}\neq0$ in general):
\begin{align}
    \begin{split}
        &\bigl[(\lambda_1+\lambda_2+\mu)z^2-\mu z-\lambda_1 z^3-\lambda_2 z^3\bigr]P(z,z)
         +\theta_1 z^2(z-1)\,\frac{\partial P}{\partial x}(z,z) \\
        &\quad =\; -\mu z(1-z)\,\pi(0,0).
    \end{split}
    \label{eq:C2:diag}
\end{align}
The unknown $\partial_x P(z,z)$ prevents a direct closed form for $P(z,z)$,
and hence for $\mathbb{E}[N]$.  In contrast to Models~A and~B$_2$, the total
mean queue length is strictly less than $\rho^2/(1-\rho)$: abandonment removes
customers from the system, reducing both $\mathbb{E}[N_1]$ and $\mathbb{E}[N_2]$
compared to the Model~A baseline.

\subsubsection*{Individual means}

The means are extracted from $P(x,y)$ by differentiation:
\begin{align}
    \begin{split}
        \mathbb{E}[N_1] \;=\; \lim_{(x,y)\to(1,1)}\frac{\partial P}{\partial x}(x,y), \qquad
        \mathbb{E}[N_2] \;=\; \lim_{(x,y)\to(1,1)}\frac{\partial P}{\partial y}(x,y).
    \end{split}
    \label{eq:C2:EN_def}
\end{align}

\subsubsection*{Little's Law with abandonment}

The relation $\mathbb{E}[N_i]=\lambda_i\,\mathbb{E}[W_i]$ holds for any
stable system.  For Model~C$_2$, $\lambda_i$ is the arrival rate of
\emph{all} class-$i$ customers, including those that will eventually abandon.
Consequently $\mathbb{E}[W_i]$ is the mean time spent in queue averaged over
both served and abandoned customers:
\begin{align}
    \begin{split}
        \mathbb{E}[W_i] \;=\; \frac{\mathbb{E}[N_i]}{\lambda_i}.
    \end{split}
    \label{eq:C2:EW}
\end{align}
The throughput $\mu(1-\pi_0)<\lambda_1+\lambda_2$ should not replace $\lambda_i$;
doing so would yield the conditional waiting time of customers who complete
service only, a different quantity.
As $\theta_1\to0^+$, the Model~A expressions~\eqref{eq:A:EN1}--\eqref{eq:A:EW2}
are recovered.

















